\newpage
\chapter{Adquisición de datos y control de instrumentos}
\vspace{3\baselineskip}

\section{Controlador del osciloscopio}
\label{sec:controlador_osciloscopio}

En esta sección se expone la implementación controlador de \emph{hardware}, perteneciente a la capa \enquote{Modelo} en la arquitectura adoptada, cuya función general es abstraer las operaciones de bajo nivel, permitiendo desacoplar la comunicación con el dispositivo del resto del sistema.

\subsection{Interfaz de comunicación y capa de abstracción de hardware}
\label{subsec:interfaz_fisica_hal}

A nivel de \emph{software}, para evitar el acoplamiento directo del código de la aplicación con las llamadas al sistema de bajo nivel (capa física de transmisión), se implementa una Capa de Abstracción de Hardware (HAL, \emph{Hardware Abstraction Layer}), bajo la Arquitectura VISA. Esta capa de abstracción permite que los comandos de consulta (\textit{queries}) y configuración (\textit{writes}) permanezcan independientes de la interfaz de hardware utilizada (sea USB, GPIB, Ethernet o RS-232), permitiendo actualizaciones de la instrumentación sin necesidad de modificar el código del presentador.

El controlador inicializa el administrador de recursos, y ejecuta un escaneo automático para buscar algún dispositivo conectado, que tenga un descriptor de recurso lógico VISA, como el siguiente:
\begin{align*}
    \texttt{ASRL1::INSTR}
\end{align*}

Si encuentra al menos un dispositivo direccionable, se conecta al mismo. La comunicación inicia instanciando el administrador de recursos VISA de Python (PyVISA), acompañado de la biblioteca \texttt{pyvisa-py} para asegurar un entorno multiplataforma sin depender de bibliotecas nativas propietarias.

Sobre el canal de transporte establecido con VISA, el software envía comandos SCPI, que debe cumplir con el siguiente formato indicado en el manual del fabricante \cite{gds1000au_ProgrammingManual}:
\begin{itemize}
    \item \textbf{No sensibilidad a mayúsculas:} Puede usarse mayúsculas y minúsculas.
    \item \textbf{Separación de parámetros:} Al configurar parámetros, debe separarse el comando y su argumento numérico, mediante un carácter de espacio en blanco simple.
    \item \textbf{Finalizador de mensaje:} Cada instrucción enviada al osciloscopio debe finalizar con el carácter delimitador de fin de mensaje \emph{Line Feed} (\textbf{\texttt{LF}}, en hexadecimal \textbf{\texttt{0x0A}} o a la secuencia de escape \textbf{\texttt{\textbackslash n}}), para que el instrumento proceda a interpretar y ejecutar la orden.
\end{itemize}

Para gestionar la configuración del osciloscopio de manera robusta, se implementa una arquitectura, basada en métodos de consulta (\emph{getters}) y configuración (\emph{setters}), que encapsulan los comandos SCPI.

Los métodos de consulta usan la instrucción \textbf{\texttt{.query()}}, reciben una cadena de caracteres, que deben convertir al tipo de dato correspondiente (\texttt{float}, \texttt{int} o \texttt{string}) y la retornan al presentador. En tanto que, los métodos de configuración emplean la instrucción \textbf{\texttt{.write()}}, sin esperar ninguna respuesta del instrumento.

Este flujo se protege con una rutina de control de excepciones que, ante fallos críticos en el bus o pérdidas de sincronismo, aborta la operación y libera los recursos del sistema de forma segura, ya que de lo contrario impediría la reconexión posterior del instrumento, debiendo reiniciar la computadora o el osciloscopio.

A continuación, se ejemplifican los métodos de consulta y configuración de la escala vertical de un canal.
\begin{python}
# Método de consulta.
def get_channel_scale(channel):
    scale = dso.query(f':channel{channel}:scale?')
    v_scale = float(scale)
    return v_scale

# Método de configuración.
def set_channel_scale(channel, value):
    dso.write(f':channel{channel}:scale {value}')
\end{python}

Por otro lado, a nivel físico, la conexión se realiza a través de un puerto USB 2.0 tipo B (esclavo) ubicado en el panel posterior del instrumento, y un puerto USB tipo A (maestro) de la computadora. Aunque, si bien el equipo se conecta a través de un puerto USB, la computadora lo detecta como un dispositivo de comunicación serie virtual\footnote{El controlador (\emph{driver}) USB, está disponible en la página del fabricante \cite{gds1000au_driver}.}, bajo la Clase de Dispositivo de Comunicación con Modelo de Control Abstracto (CDC-ACM, \emph{Communication Device Class - Abstract Control Model}), que aparece en el administrador de dispositivos con el siguiente nombre \cite{gds1000au_ProgrammingManual}:
\begin{align*}
    \texttt{COM1}
\end{align*}
Donde \texttt{COM1} identifica el número de puerto serie virtual asignado por el sistema operativo.

\subsection{Sistema de Adquisición y Almacenamiento}
\label{subsec:arquitectura_flujo}

La adquisición del impulso de alta tensión desde el generador Marx, hasta el almacenamiento en la computadora, se estructuran de la siguiente forma:

\begin{enumerate}
    \item Acondicionamiento analógico.
    \item Captura de impulso
    \item Digitalización.
    \item Desempaquetado de datos.
    \item Reconstrucción y escalamiento de la onda.
\end{enumerate}

\subsubsection{Acondicionamiento analógico}

La señal eléctrica, debe ser adaptada en amplitud antes de ingresar al osciloscopio, para no dañarlo.

Para el canal de tensión (CH1), la señal se deriva utilizando un divisor de tensión resistivo, acoplado en cascada serie con un atenuador coaxial, por lo que el factor de atenuación equivalente resulta:
\begin{equation}
    \label{eq:atenuacion_tension}
    A_{\text{CH1}} = \text{Divisor} \times \text{Atenuador}
\end{equation}

Para el canal de corriente (CH2), la descarga atraviesa una resistencia shunt conectada en serie con el elemento bajo ensayo, de la cual se toma su tensión a bornes, acoplada a un atenuador que reduce su magnitud a niveles adecuados.
\begin{equation}
    \label{eq:atenuacion_corriente}
    A_{\text{CH2}} = R_\text{Shunt} \times \text{Atenuador}
\end{equation}

Ambas señales llegan al instrumento por medio de cables coaxiales de \qty{50}{\ohm}, para reducir acoplamientos electromagnéticos parásitos (EMI).

\subsubsection{Captura de impulso}
\label{subsec:captura_impulso}

La captura de impulsos (\qty[parse-numbers=false]{1.2/50}{\micro\second}) requiere que el osciloscopio opere en modo de disparo único (\emph{Single Trigger}), dado que es un transitorio no periódico.

Cuando el operador inicia el modo de captura desde la GUI, el presentador crea y ejecuta un subproceso asíncrono, que mediante el controlador de \emph{hardware}, realiza un sondeo periódico (\emph{polling}) al osciloscopio, para saber el estado de disparo del mismo.

Al detectar el impulso finaliza el bucle de \emph{polling} y emite una señal, que es capturada en el hilo de ejecución principal por un receptor (\emph{slot}) del presentador, y finaliza el subproceso.

Como el instante de descarga es indeterminado, se implementa un hilo de ejecución secundario asíncrono, dedicado a consultar el estado de disparo, ya que si se ejecutara en el hilo principal, sería una operación bloqueante, que congele la aplicación e inhabilite la GUI. Por lo que, así se mantiene siempre reactivo el hilo principal de eventos gráficos, permitiendo al usuario manipular oscilogramas previos, cargar campos o cualquier otra acción, mientras el sistema espera la captura del impulso.

\subsubsection{Digitalización}

Las señales eléctricas que ingresan a los bornes analógicos del osciloscopio, se muestrean a una tasa que depende de la base de tiempo utilizada. Después, el convertidor analógico-digital (ADC) realiza la cuantización en amplitud, dividiendo el rango dinámico vertical en $2^8 = 256$ niveles de tensión discretos (\num{8} bits). A continuación, se graban en la memoria del instrumento.

\subsubsection{Desempaquetado de datos}
\label{subsec:desempaquetado_datos}

Luego de que el osciloscopio digitalice la señal, el presentador ordena la descarga de los datos brutos (\emph{RAW}) de los canales habilitados, y recibe un bloque de datos, compuestos por el encabezado y el vector de datos de la onda, que cumple con la siguiente estructura\footnote{La estructura de datos descrita en el manual del fabricante tiene algunas diferencias con los resultados obtenidos del instrumento, debido a la dificultad que esto presentó al momento de intentar procesarla, se contactó con el frabricante obteniendo un script de procesamiento del paquete de datos, que se corresponde con el formato indicado en este trabajo. A su vez, se encontró que el manual de programación de la series GDS-1000B aporta detalles explicativos sobre cómo procesar esta estructura \cite{gds1000b_ProgrammingManual}} \cite{gds1000au_ProgrammingManual}:

\begin{equation*}\label{eq:formato_memoria}
\text{\texttt{\# A B C D E}}
\end{equation*}
\begin{enumerate}
    \item[\texttt{\#:}] [\num{1} byte] Caracter ASCII que indica el inicio de la cabecera.
    \item[\texttt{A:}] [\num{1} byte] Dígito ASCII que especifica la cantidad de caracteres que componen el campo \texttt{B}. Para adquisiciones estándar toma el valor ASCII \texttt{'4'}, mientras que para adquisiciones de memoria extendida toma el valor \texttt{'7'}.
    \item[\texttt{B:}] [\num{4} o \num{7} bytes] Cadena de caracteres ASCII que representa el tamaño total en bytes del bloque de datos que le sigue (\texttt{C}+\texttt{D}+\texttt{E}).
    Para adquisiciones estándar toma el valor ASCII \texttt{'8008'}, mientras que para adquisiciones de memoria extendida puede tomar los valores \texttt{'2000008'} o \texttt{'4000008'}.
    \item[\texttt{C:}] [\num{4} bytes] Intervalo de tiempo ($\Delta t$) entre dos muestras consecutivas de la señal, en formato de punto flotante de precisión simple según IEEE-754 (\num{32} bits), en orden \emph{little-endian} (se envía primero el byte menos significativo):
    \begin{equation*}\label{eq:periodo_muestreo}
    \Delta t = t[n+1] - t[n]
    \end{equation*}
    \item[\texttt{D:}] [\num{4} bytes] Bloque de datos no utilizado.
    \item[\texttt{E:}] [\num{8e3}, \num{2e6} o \num{4e6} bytes] Vector de datos de la onda, cuya longitud es definida en el campo \texttt{B}. Cada valor entero (base \num{10}) es expresado con \num{2} bytes (\num{16} bits), y transmitidos en orden \emph{big-endian} (se envía primero el byte más significativo).
\end{enumerate}

Una vez finalizada la descarga de datos, el presentador ordena al controlador de \emph{hardware} que decodifique cada bloque acorde a su formato, considerando la estructura recién mencionada.

\subsubsection{Reconstrucción y escalamiento de la onda}

Con el vector de valores enteros brutos decodificado, primero se aplica la ecuación \autoref{eq:reconstruccion_tension} para reconstruir la forma de onda que llegó a bornes del osciloscopio, considerando la escala vertical configurada al momento de la captura, y el Factor AD = \num{25} definido en el manual de programación del fabricante \cite{gds1000au_ProgrammingManual}:
\begin{equation}
\label{eq:reconstruccion_tension}
\text{Onda sin escalar}[n] = \text{Datos brutos}[n] \times \frac{{\text{V/div}}}{\text{Factor AD}}
\end{equation}
A continuación, para obtener las magnitudes reales aplicadas al objeto bajo ensayo, se multiplica escalarmente los vectores de tensión y corriente, por sus respectivos factores de atenuación, cargados por el usuario en los campos pertinentes de la ventana, aplicando la siguiente ecuación:
\begin{equation}
\label{eq:reconstruccion_tension_total}
{\text{Onda escalada}}[n] = \text{Onda sin escalar}[n] \times A_{\text{CH<X>}}
\end{equation}
Los vectores resultantes de tensión (en volts, \unit{\volt}) y de corriente (en amperes, \unit{\ampere}), se transfieren al módulo de procesamiento matemático para determinar sus parámetros normativos: tiempo de frente $T_1$, tiempo de cola $T_2$, valor de cresta $U_{\text{t}}$ y sobreimpulso OS, conforme a los requerimientos de la norma \cite{iec60060_1}.






\newpage

\noindent
===============================================================================


\noindent
===============================================================================